\documentclass[t]{beamer}
\usepackage[utf8]{inputenc}
\usepackage[francais]{babel}
\usepackage{tikz,pgf}
\usepackage{wrapfig}
\usepackage{times}
\usepackage[absolute,overlay]{textpos}
\usepackage{url}
% \usepackage[texcoord,grid,gridunit=mm,gridcolor=red!10,subgridcolor=green!10]{eso-pic}

\definecolor{bleu}{rgb}{0.282, 0.439, 0.643}

\usecolortheme[named=bleu]{structure}
\useinnertheme[shadow]{rounded}

\hypersetup{pdfpagemode=FullScreen}

\setbeamersize{text margin left=0.6cm}
\setbeamersize{text margin right=0.4cm}
\setbeamersize{text margin left=.5cm, text margin right=.5cm}
\setbeamertemplate{navigation symbols}{}

\AtBeginSection[]{
    \begin{frame}
  	 \frametitle{\vspace*{-1cm} \begin{flushright} \ \ \vskip-0.25cm \end{flushright} }
        \vspace*{-0.75cm}
	    \begin{center}
            \centering
            \tableofcontents[currentsection]
	    \end{center}
    \end{frame}
}

\begin{document}
    \begin{frame}
        \pgfputat{\pgfxy(-0.505,0.685)}{\pgfbox[left,top]{\pgfimage[width=\paperwidth,height=\paperheight ]{./assets/UTT_garde_empty.pdf}}}
        \vspace*{2.5cm}

        \begin{center}
            \textcolor{white}{
                {\large Soutenance de TX} \\
                {\footnotesize   Printemps 2015}\\
            \vspace*{0.5cm}
                {\footnotesize \textbf{Développement d'une de fonctions de chiffrement en C}}\\
            \vspace*{0.25cm}
                {\footnotesize Pergoud Florent\\ Labate Aurélien}\\
          }
        \end{center}
    \end{frame}

\setbeamertemplate{background canvas}{
    \includegraphics[width=\paperwidth,height=\paperheight]{./assets/UTT_fond.pdf}}
    \addtobeamertemplate{footline}{
    \vspace{-0.5cm}
    \begin{flushright}
    \insertframenumber/\inserttotalframenumber
    \end{flushright}
}


\section{Introduction}
    \begin{frame}[c]
        \begin{center}
            \frametitle{\vspace*{-1.0cm} \begin{flushright} Introduction \ \ \vskip-0.25cm \end{flushright} }
            \vspace*{-0.75cm}
            \begin{itemize}
                \item Acquérir des connaissances en cryptographie 
                \item Création d'une librairie en C avec des fonctions de cryptographie
                \item Faire un programme d'exemple avec un menu pour présenter les algorithmes
            \end{itemize}
        \end{center}
    \end{frame}


\section{Environnement de programmation}
    \begin{frame}[c]
        \frametitle{\vspace*{-1.0cm} \begin{flushright} Environnement de programmation \ \ \vskip-0.25cm \end{flushright} }
        \vspace*{1.5cm}
        \begin{center}
            \textbf{Une architecture modulaire ?}
        \end{center}
        Inconveniants
        \begin{center}
            \begin{itemize}
                \item Plus cher
                \item Plus volumineux
                \item Plus de cables
            \end{itemize}
        \end{center}
        Avantages
        \begin{center}
            \begin{itemize}
                \item Réutilisable
                \item Facilité de séparation des tâches
                \item Changement facile de chaques modules
                \item Moins de problèmes en cas d'erreurs
            \end{itemize}
        \end{center}
    \end{frame}
    \begin{frame}[c]
        \frametitle{\vspace*{-1.0cm} \begin{flushright} Architecture modulaire \ \ \vskip-0.25cm \end{flushright} }
        \vspace*{1.5cm}
        \begin{center}
            \textbf{Le coeur}
        \end{center}
        \begin{center}
            \begin{itemize}
                \item Commande tous les autres modules
                \item Plus de puissance de calcul
                \item Language plus haut niveau
            \end{itemize}
        \end{center}

        Les autres modules :
        \begin{center}
            \begin{itemize}
                \item Moteur
                \item Pince
                \item Capteurs
            \end{itemize}
        \end{center}
    \end{frame}
    \begin{frame}[c]
        \frametitle{\vspace*{-1.0cm} \begin{flushright} Architecture modulaire \ \ \vskip-0.25cm \end{flushright} }
        \vspace*{1.5cm}
        \begin{center}
            \textbf{Module moteur}
        \end{center}
        \begin{center}
            \begin{itemize}
                \item Récupère les informations des roues codeuses
                \item Commande et asservis les moteurs
                \item Calcul la position et la trajectoire
            \end{itemize}
        \end{center}
    \end{frame}
    \begin{frame}[c]
        \frametitle{\vspace*{-1.0cm} \begin{flushright} Architecture modulaire \ \ \vskip-0.25cm \end{flushright} }
        \vspace*{1.5cm}
        \begin{center}
            \textbf{Module pince}
        \end{center}
        \begin{center}
            \begin{itemize}
                \item Génère les impulsions à fréquence optimale pour controller les moteurs
                \item Initialise la position des moteurs pas à pas
            \end{itemize}
        \end{center}
    \end{frame}
\section{Les algorithmes principaux}
        \begin{frame}[c]
            \frametitle{\vspace*{-1.0cm} \begin{flushright} Les algorithmes principaux \ \ \vskip-0.25cm \end{flushright} }
            \vspace*{0.5cm}
			\begin{center}
			\textbf {Les différents algorithmes}
			\end{center}
            \begin{center}
                \begin{itemize}
                    \item Exponentiation rapide
                    \item Test de primalité de Rabin-Miller
                    \item Diffie-Helmann
                    \item ElGamal
                    \item RSA
                \end{itemize}
            \end{center}
        \end{frame}

        \begin{frame}[c]
            \frametitle{\vspace*{-1.0cm} \begin{flushright} Les algorithmes principaux \ \ \vskip-0.25cm \end{flushright} }
            \vspace*{0.5cm}
             \begin{center}
             \textbf{Test de primalité de Rabin-Miller}
             \end{center}
            \begin{center}
                \begin{itemize}
                    \item Algorithme probabiliste
                    \item En entrée un entier n à tester et un entier k (précision)
                    \item En sortie "probablament premier" ou "composé"
                    \item Probabilité de se tromper :  $4^{-k}$ 
                \end{itemize}
            \end{center}
        \end{frame}

        \begin{frame}[c]
            \frametitle{\vspace*{-1.0cm} \begin{flushright} Les algorithmes principaux \ \ \vskip-0.25cm \end{flushright} }
            \vspace*{0.5cm}
             \begin{center}
             \textbf{ElGamal}
             \end{center}
            \begin{center}
                \begin{itemize}
                    \item Algorithme servant à chiffrer et déchiffrer un message 
                    \item Clé publique d'Alice : (p,g,A)
                    \item Clés privées a et b avec $A=g^a$ et $B=g^b$
                    \item Message chiffré de Bob : (M,B)
                    \item Déchiffrement $m=M/B^a$
                \end{itemize}
            \end{center}
        \end{frame}

        \begin{frame}[c]
            \frametitle{\vspace*{-1.0cm} \begin{flushright} Les algorithmes principaux \ \ \vskip-0.25cm \end{flushright} }
            \vspace*{0.5cm}
             \begin{center}
             \textbf{RSA}
             \end{center}
            \begin{center}
                \begin{itemize}
 				  \item Algorithme servant à chiffrer et déchiffrer un message 
                    \item Génération d'un couple de clés : clé publique (n,e), clé privée (n,d)
                    \item Généralement on prend un e constant
                    \item $n=p\times q$ ou p et q sont premiers
                    \item $d=e^{-1} \mod\phi(n)$ où $\phi(n) = (p-1)(q-1)$
                    \item chiffrement : $C=M^e \mod(n)$
                    \item déchiffrement : $M=C^d \mod(n)$
                \end{itemize}
            \end{center}
        \end{frame}

\section{Conclusion}
    \begin{frame}[c]
        \frametitle{\vspace*{-1.0cm} \begin{flushright} Conclusion \ \ \vskip-0.25cm \end{flushright} }
        \vspace*{-0.75cm}
        \begin{itemize}
            \item Beaucoup de connaissances acquises ou approfondies:
            \begin{itemize}
            		\item Notions et principes de crytographie
            		\item C++
            		\item GMP
            		\item \LaTeX
            		\item Gestion de projet de recherche
			\end{itemize}
             
            \item La bibliothèque a été crée et est fonctionnelle
            \item Beaucoup d'exemples pour utiliser la bibliothèque et comprendre les algorithmes de cryptographie en jeu
            
        \end{itemize}

    \end{frame}



\end{document}